% A user-declared pgfmath function returns its result in the form its body leaves it (an integer
% `\pgfmathresult` stays an integer) and receives an integer-form argument as an integer, so a
% `?:`, `ifthenelse` or `\ifnum` on its result behaves as in pgf.
% witness: none (probe; review of 59b).
% status:  GREEN (59b)
% oracle:  lualatex 0 errors: "[1][1][1.0][7][6.0]YES".
% expect:  0 errors, 0 warnings.
\documentclass{article}
\usepackage{pgf}
\pgfmathdeclarefunction{pos}{1}{\pgfmathparse{#1>0}}
\pgfmathdeclarefunction{one}{0}{\def\pgfmathresult{1}}
\pgfmathdeclarefunction{idf}{1}{\edef\pgfmathresult{#1}}
\begin{document}
\pgfmathparse{1>0?pos(3):0}[\pgfmathresult]\pgfmathparse{1>0?one:0}[\pgfmathresult]%
\pgfmathparse{one+0}[\pgfmathresult]\pgfmathparse{idf(7)}[\pgfmathresult]%
\pgfmathparse{idf(2*3)}[\pgfmathresult]%
\pgfmathsetmacro\n{pos(3)}\ifnum\n=1 YES\else NO\fi
\end{document}
